% !TeX root = ../papers/nc-p3.tex
\section{非被約な四次完全交叉上の全鉛筆}
\label{sec:nonreduced-quartic-pencils}
\begin{theorem}\label{thm:nonreduced-quartic-pencils}
$C\subset\mathbb P^3$を非被約な二二次式完全交叉、$L,M$を各既約成分で正次数、
全次数4の可逆層、$W\subset H^0(C,LM)$を基点自由鉛筆とする。
反転二窓は全て標準Koszul AS正則である。$L=M$（同型を略記）では対称完全交叉となり、
$L\ne M$では反対階数8で、添字保持同型は$(C,L,M,W)$の順序付き曲線同型と同値である。
特に後者では反対核から被約台とその冪零構造も決まる。
\end{theorem}

\begin{proof}
まず埋込みを保つ乗法群で$H^1(\mathcal O_C)$に非自明に作用するものを作る。
二式は共通因子を持たず、一次元Cohen--Macaulay性により重複する既約成分が存在する。
この成分は次数から直線か滑らかな平面二次曲線である。後者のイデアルを$(h,f)$と書くと、
二式のうち一方は$h$の倍数にできる。二次曲線の一般点で交叉が単純にならないため
その倍数は$h^2$である（他方は$f$を非零係数で含む）。
他方の式の$h$一次項は非退化三変数二次式の変数の平行移動で除ける。

重複直線のイデアルを$(x,y)$とする。二式の$(x,y)^2$を法とした係数の二行二列行列は、
直線上で恒等的に行列式零である。階数1行列の直線はSegre曲面の母線なので、
共通左核または共通右核を持つ（射影像が一点以下の場合も含む）。
左核があれば鉛筆内に$x,y$の二次式があり、右核があれば線型部分を$x\ell_i(z,t)$にできる。
$\ell_i$が従属なら再び二変数の式を得る。独立なら$z,t$に$x,y$の倍数を加えて
$xz+a y^2,xt+b y^2$となり、$y^2$を消す線型結合が二平面を与える。
従っていずれも鉛筆に重平面または二平面がある。

重平面$h^2$を取ると、残る式の平面への制限の階数に応じて、滑らかな二次曲線のほかは
$(h^2,xy+\kappa hz)$または$(h^2,x^2+\kappa hz)$である
（$\kappa=0,1$、不要な変数を補って四変数とする）。
相異なる二平面$xy$を取る場合、その交線上に残る式が恒等零でなければ、
非被約性から片方の平面切断は二重直線である。そこでその制限を$z^2$として、
残りを$x$と一次式の積で書く。$z$に$x$の倍数を加え、$xy$を除くと、
第二式は$z^2+xt,z^2+x^2,z^2$のいずれかとなる。
交線全体が含まれれば、第二式は$x(\ell+a x)+y(m+b y)$と書ける。
交線上の$\ell,m$が独立なら$xz+yt$、非零比例なら$z(x+y)$にできる。
一方だけ非零なら共通因子を避けるため$xz+y^2$にできる。
両方零なら二変数正則列であり、二変数の平方の円錐曲線と鉛筆の二交点から
$(x^2,y^2)$にできる（接線なら共通一次因子が生じる）。これらの表示の重複は許す。

二式の重みを$e_1,e_2$、座標の重みを$w_i$とすると、随伴公式から
$H^1(\mathcal O_C)$の指標は符号を除いて$\sum w_i-e_1-e_2$である。
未使用変数があればそれだけに重み1を置く。滑らかな二次曲線の重複では$h$だけに1を置く。
残る$(h^2,xy+hz)$では$(h,x,y,z)$に$(0,1,1,2)$を、
$(xy,z^2+xt),(xy,xz+yt)$では$(x,y,z,t)$にそれぞれ$(0,0,1,2),(0,0,1,1)$を置く。
全て非零指標となる。

この表示から被約台は滑らかな有理曲線、二成分が一点で横断する木、三直線の鎖、
または三本の独立な方向を持つ同一点の三直線である。三直線の最後の場合も
正規化からの降下条件は値の一致だけで、全ての場合に被約台の$\operatorname{Pic}^0$は零である。
各成分の次数に一般点の長さを掛けて足すと4なので、正の多重次数は埋込み束
$L_0=\mathcal O_C(1)$のものに一意に決まる。実際、重複二次曲線は次数2、
四重直線は次数1、二重直線と単純二次曲線なら次数$(1,2)$で、
他の複数直線の場合は重複度の和が4なので全て次数1である。
単元列と冪零イデアルの対数から、多重次数零の可逆層全体は加法群である。
その接空間は$H^1(\mathcal O_C)=k$なので$\operatorname{Pic}^0(C)=\mathbb G_a$となる。

重複成分$P\simeq\mathbb P^1$に対し
$I=\mathcal Hom_C(\mathcal O_P,\mathcal O_C)\subset\mathcal O_C$を取る。
有限閉埋込みの双対性と$\omega_C=\mathcal O_C$により$I=\omega_P$で、
$H^1(I)\to H^1(\mathcal O_C)$は定数の制限の双対だから同型である。
$I$は一般点でArtin局所環のsocleにあり、$P$への像は零である。
従って$I^2=0$である。$\mathcal O_C\to\mathcal O_P\xrightarrow{d}\omega_P=I$の合成は
平方零の導分$\delta$で、$g_s=1+s\delta$は自己同型を与える。
$1+I$からの写像$H^1(I)\to\operatorname{Pic}^0(C)$は微分が同型の加法群の準同型であり、
標数零だから同型である。遷移関数$f$の変化は$1+s\,d\log(f|_P)$なので、
$g_s^*$は$P$上の次数$d$の束を、そのPicard座標で$sd$だけ移す（座標の定数倍を許す）。
従って任意の正または負の冪$L_0^n$の多重次数の束は全てその自己同型像である。
また$H^0(\mathcal O_C/I)=k$は上の$H^1$の同型から従う。多重次数零の非自明な束は
非零の$1+I$コサイクルで貼られ、商上の定数1は持ち上がらず、$H^0(I)=0$なので切断は零である。

第一束を$L=L_0$に揃え、第二束を$M_s=g_s^*L$とする。
上で構成した乗法群は$L$を固定し、$\operatorname{Pic}^0(C)$に非零指標で作用するので、
$s\ne0$の順序付き束の対は一軌道である。
交互切断環（偶奇二頂点に折り畳む）の各長さ$n>0$の束は、
上の平行移動により$L^n$と同じコホモロジーを持つ。
埋込み完全交叉の分解から$h^0=4n,h^1=0$で、負の冪には切断がない。
従って切断空間はパラメータ直線上で束となり基底変換と可換である。
$s=0$の切断環は二次完全交叉の斉次座標環を二頂点に引き戻したものでKoszulである
（二次正則列で多項式環を順に割る）。内部次数を一つ固定した相対正規化bar複体は
有限長のベクトル束の複体である。非対角Torの消滅はこの次数ごとに原点の開近傍へ移る。
各近傍が非零軌道に会うので、各次数で全パラメータ上の消滅を得る。
特に任意のこの多重次数の二束の切断の積は全射である。

交互切断環から$W$の二つの中心二次元で割る。一般の$a\in W$は全既約成分の一般点で非零、
ゆえに正則元である。第二の生成元は$a$の零因子上で単元である。
従って切断の等式$bu=av$から$u$は$a$で割れ、負次数部分の切断も零なので、
二元は正則列になる。定理\ref{thm:elliptic-pencil-all-as}の鉛筆完全列とSerre双対性により、
商はHilbert列$1,4,6,4,1$のKoszul Frobenius代数で、その二次双対が求める標準AS正則代数である。
同じ束なら二つの定義式と鉛筆が四二次式の完全交叉を与える。
異なる束なら$L/M$とその逆の切断が零で、同じ鉛筆列は$H^0(L^2)$と$H^0(M^2)$に
完全対合を与える。両束は$L_0$の自己同型像だから二次核は両埋込みの定義式そのものである。

最後に両埋込みの間の同定をスキーム上で一意に戻す。
階数8の両反対核には同時に四つの基底を二つの射影基底に還元する既出の議論が適用できる
（定理\ref{thm:elliptic-pencil-reconstruction}の証明）。
第一の候補グラフで両曲線を$C$と同定し、第二候補を自己同型$g$のグラフとする。
$s,t\in H^0(L)$の行列式$s(x)t(y)-t(x)s(y)$に、
任意の$H^0(L)\boxtimes H^0(M^2/L)$の切断を掛ける。積は両側の二次切断のテンソルで、
第一グラフで零だから与えられたテンソル汎関数で消える。
第二グラフでは$D=L\,g^*L$, $N'=L\,g^*M$として、
乗じた切断の積が$H^0(N'^2/D)$を張る。実際、二因子$L,g^*(M^2/L)$の
多重次数はともに一意な正の型なので先ほどの積の全射性が使える。
第二候補の鉛筆traceにおいて$H^0(D)$と$H^0(N'^2/D)$は完全に対合する。
これは$D/N'=g^*(L/M)$とその逆の切断が零だからである。
よって$s\,g^*t-t\,g^*s=0$が切断として成り立つ。$|L|$は閉埋込みなので、
切断の比を局所座標に使えば$g$は冪零方向を含め恒等である。
鉛筆は隣接核を積で写した像として戻り、逆に順序付きデータの同型は全ての積を保つ。
\end{proof}
